\section{时间的刚性：不可压缩的过程}

\subsection{物理过程的非零时长属性}

资源的账本回答的是物质从哪里来，还没有回答它要多久才能就位。加速不减少任何载荷，它只改变载荷通过系统的速率。同一份材料、同一份能量、同一份热和同一份误差，被要求在同一条更短的窗口里通过系统。当时间成为约束，空间的三种边界不再分头出现，而是转换成同一个量：单位时间内的负荷。

把时间当作可以压缩的资源，是智能时代最普遍的假设。这个假设在信息侧确实成立：等待可以被消除，重试可以被并行，搜索可以被剪枝，一个环节的延迟不会自动传递给下一个环节。但在物质侧，同一个状态转移，占用的材料、注入的能量、产生的热、积累的损伤和维护需求都要在更短的窗口里被处理完。

AI时代最容易产生的错觉，是把速度等同于效率。模型可以快速生成答案，系统可以快速响应信号，平台可以快速分发内容，组织可以快速获得指标，于是人们很自然地认为越快越好，只要把过程压缩，把等待消除，把反馈拉高频，现实系统就会变得更有效率。但时间不是一个可以任意删除的空白。在现实系统中，一个状态转移总要穿过材料、能量、设备、身体、组织和环境，把同一个状态转移压进更短时间，并不只是“少花时间”，而是提高单位时间内穿过系统的通量。时间越短，系统必须在同一窗口内处理更多能量、更多物质、更多热、更多误差和更多风险，加速的本质，不是速度本身，而是通量密度上升。

反应、扩散、养护、生长与服役各有时间尺度。计算可以帮助改变配方、设备、温度、顺序与并行方式，从而缩短其中一些过程；某些长期结果仍须通过足够暴露才能确认。设计速度、制造速度与获得证据的速度，需要分别记录。

\subhead{通量密度}

通量密度可以用一个简单的量描述：
\begin{equation}
J=\frac{\Delta X}{A\,\Delta t},
\label{eq:flux}
\end{equation}
其中$J$为通量密度，表示单位截面上单位时间内通过的载荷量；$\Delta X$为这次状态转移需要穿过的载荷总量，按对象不同可以是能量、物质、热量、误差或需要承接的后果，量纲随对象而定；$A$为系统的通量截面，即系统在单位时间内能够承受该载荷的有效截面积；$\Delta t$为完成这次转移所允许的时间窗口，单位秒。这个关系成立的前提是系统处在同一工作区间内，即截面$A$不随负荷提高而自动扩大。它给出的是一个比例判断：在$\Delta X$和$A$不变时，把$\Delta t$缩小到原来的十分之一，$J$就上升到十倍，而系统能否承受十倍的通量密度，与$J$的升高无关，只由$A$决定。同样一件事，用十天完成和用一天完成，系统面对的不是同一个过程：能量要更集中地注入，废热要更快排出，设备要承受更高负荷，人员要减少修正时间，错误会更快扩散，失败也更少有机会被早期截断。时间被压缩之后，并不会消失，而会以功率、热流、应力、检测压力和修复压力的形式重新出现。

因此，加速能否成立，不取决于目标时间能否写进计划，而取决于系统是否有足够的通量截面。所谓通量截面，是系统承受单位时间流量的能力，它包括供能截面、散热截面、承载截面、检测截面、修复截面和责任截面。供能截面不足，过程无法驱动；散热截面不足，系统会升温；承载截面不足，结构会疲劳；检测截面不足，缺陷会穿过系统；修复截面不足，小错误会被复制成大故障；责任截面不足，快速行动后的后果无人承接。这比“速度有代价”更根本，问题不是快本身，而是快把原本分散在时间里的负荷压缩到了更窄的窗口。时间原本承担了一部分缓冲功能，它让热有时间扩散，让材料有时间释放应力，让人员有时间检查，让组织有时间协调，让错误有时间显影。加速切掉这些缓冲，就必须用更强的结构、更高的冗余、更快的检测和更大的维护能力来替代。

这里有一种常见的辩解：既然缓冲可以用技术替代，加速的极限就只是投入多少的问题，而不是原理问题。这个说法在单个环节上常常成立，放到系统上却不成立，因为六类截面并不彼此独立。加厚结构能提高承载截面，同时增加质量和热容；提高刷新率能扩大检测截面，同时增加能耗和误报；增加冗余能提高修复截面，同时提高协调成本；把决策交给自动过程能加快责任截面的通过速度，同时削弱责任归属。加宽一个截面往往压缩另一个截面，所以加速不是一次可以买断的投资，而是必须持续维持的结构代价。这也是为什么加速不是线性的：把十天压缩到五天，可能只是减少等待、优化排程、削掉返工；把五天压缩到一天，可能就需要更多设备、更强并行、更高库存、更密集检查；把一天压缩到一小时，可能已经触及散热、可靠性、安全边界和责任承载。越接近系统极限，时间每缩短一分，后台强度往往不止增加一分，边际加速阻抗会迅速上升。

混凝土的强度随水化发展，28天是常用检验龄期，实际养护要求则随材料和工程条件确定。大体积浇筑还要控制水化热与温差，施工分层、冷却和保温都在改变可行的节奏。锂离子电池的高倍率充电则可能受到传输、嵌入和析锂等机制限制。加速这些过程，需要改变材料、温度、结构或操作方式，再检验相应的成本与寿命。

通信和计算提供了另外两个界。光纤中的光速只有真空中的约三分之二，跨洋通信的单程延迟因此在几十毫秒的量级；增加波长、提高功率、加密中继都能扩大带宽，却不能让信号提前到达，带宽是可以扩的截面，延迟不是。单核时钟频率在 2005 年前后停在了数 GHz 量级，继续提高频率会让功耗密度超过散热能力，算力的增长因此转向多核并行。并行能够摊销的只是彼此独立的步子，不能减少一条依赖链上的步数，这条界线在后面的讨论中还要反复用到。

\subsection{功率与时间的互斥方程}

通量密度上升还只是一个描述，没有说明负荷落到哪些量上。把同一件事压进更短的时间窗口，时间挤出来的负荷会依次出现在几个可以分别考察的地方，它们的次序不是随意的，前一个量决定了后一个量的规模。
\begin{figure}[ht]
\centering
% >>>>>>>>>> figs/fig-power-time.tex >>>>>>>>>>
% 图 2-3  功率与时间的互斥，以及压不动的功率天花板
% 双对数坐标：1 个数量级 = 横向 1.25 cm、纵向 0.72 cm。
% 固定能量需求 E 对应等能线 P·t = E；沿同一条线只能「用功率换时间」，
% 整条线下移一格（能量需求降一个数量级）才叫效率改善，斜率不变。
% 水平虚线是材料强度、传热密度、电网容量共同设的功率天花板；
% 它与等能线的交点给出最短时间 t_min——这才是时间压不下去的原因。
\begin{tikzpicture}[x=1cm, y=1cm, font=\small,
  ax/.style={-{Stealth[length=2.0mm]}, ink, line width=0.6pt},
  tk/.style={font=\scriptsize, text=slate, anchor=north},
  yk/.style={font=\scriptsize, text=slate, anchor=east}
]
% 上缘原先 1.30，纵轴的 $E$ 标注与「：现行路线的能量需求」距页边只剩约 2pt。放宽到 1.55。
\useasboundingbox (-1.35,1.55) rectangle (12.75,-6.30);

\def\aX{1.05} \def\aY{-4.90} \def\dx{1.25} \def\dy{0.72} \def\aR{11.15}
\pgfmathsetmacro\yTop{\aY+\dy*9.2}

\draw[ax] (\aX,\aY) -- (\aR,\aY) node[anchor=west, font=\scriptsize, text=ink] {时间 $t$ (s)};
\draw[ax] (\aX,\aY) -- (\aX,\yTop)
      node[anchor=south, font=\scriptsize, text=ink, yshift=1pt] {功率 $P$ (W)};

\foreach \e in {0,...,7} {
  \pgfmathsetmacro\xx{\aX+\dx*\e}
  \draw[slate!55, line width=0.4pt] (\xx,\aY) -- (\xx,\aY-0.07);
  \node[tk] at (\xx,\aY-0.10) {$10^{\e}$};
}
\foreach \e in {0,2,4,6,8} {
  \pgfmathsetmacro\yy{\aY+\dy*\e}
  \draw[slate!55, line width=0.4pt] (\aX,\yy) -- ({\aX-0.07},\yy);
  \node[yk] at ({\aX-0.09},\yy) {$10^{\e}$};
}

% ---------- 守恒律禁区：最下面那条等能线之下（先铺底，线压其上）----------
\begin{scope}
\clip (\aX,\aY) rectangle (\aR,\yTop);
\fill[slate!10] ({\aX+\dx*1},{\aY+\dy*7}) -- ({\aX+\dx*7},{\aY+\dy*1})
      -- ({\aR},{\aY+\dy*1}) -- ({\aR},\aY) -- ({\aX+\dx*1},\aY) -- cycle;
\end{scope}

% ---------- 两条平行的等能线 ----------
% E=10^9：f = 9 − e；E=10^8：f = 8 − e。两线相距恰好一个数量级。
\draw[terra!90, line width=1.7pt] ({\aX+\dx*1},{\aY+\dy*8}) -- ({\aX+\dx*7},{\aY+\dy*2});
\draw[forest!85, line width=1.5pt, dashed] ({\aX+\dx*1},{\aY+\dy*7}) -- ({\aX+\dx*7},{\aY+\dy*1});
\node[font=\scriptsize, text=terra!90, anchor=south west] at ({\aX+\dx*1.02},{\aY+\dy*8.10})
      {$E=10^{9}$ J：现行路线的能量需求};
\node[font=\scriptsize, text=forest!85, anchor=south east] at ({\aX+\dx*7.0},{\aY+\dy*1.12})
      {$E=10^{8}$ J：效率改善后};
\draw[ink!55, line width=0.5pt, {Stealth[length=1.4mm]}-{Stealth[length=1.4mm]}]
      ({\aX+\dx*6.30},{\aY+\dy*2.70}) -- ({\aX+\dx*6.30},{\aY+\dy*1.70});
\node[font=\scriptsize, text=ink, anchor=west] at ({\aX+\dx*6.42},{\aY+\dy*2.20})
      {一格};

% ---------- 功率天花板：说明块放在线右侧上方空白，短箭头指回交点 ----------
\def\ceil{\aY+\dy*6.6}
\draw[rust!90, line width=1.1pt, densely dashed] ({\aX+0.05},\ceil) -- ({\aR-0.05},\ceil);
\node[font=\scriptsize, text=rust!90, anchor=south west, align=left]
      at ({\aX+\dx*3.30},{\aY+\dy*6.78})
      {功率天花板：材料强度、传热密度、电网容量\\[1pt]天花板以上不可达，任何工艺都进不去};
\draw[rust!75, line width=0.6pt, -{Stealth[length=1.6mm]}]
      ({\aX+\dx*3.18},{\aY+\dy*6.62}) -- ({\aX+\dx*2.52},{\aY+\dy*6.62});

% ---------- 交点给出最短时间 ----------
\pgfmathsetmacro\xmin{\aX+\dx*2.4}
\draw[slate!65, line width=0.5pt, dotted] (\xmin,\aY) -- (\xmin,\ceil);
\draw[ink, line width=1.6pt] (\xmin-0.12,\aY) -- (\xmin+0.12,\aY);
\node[font=\scriptsize, text=ink, anchor=north] at (\xmin,\aY-0.10) {$t_{\min}$};

% ---------- 沿线设备 ----------
\foreach \ex/\ey in {1.4/7.6, 2.4/6.6, 4.0/5.0, 6.0/3.0} {
  \node[circle, fill=ink, inner sep=1.4pt] at ({\aX+\dx*\ex},{\aY+\dy*\ey}) {};
}

% ---------- 禁区注记 ----------
\node[font=\scriptsize\heiti, text=slate, anchor=north east]
      at ({\aR-0.15},{\aY+\dy*0.58})
      {等能线之下：能量不够，任何工艺都做不到};
\end{tikzpicture}

% <<<<<<<<<< figs/fig-power-time.tex <<<<<<<<<<

\caption{功率与时间的互斥。固定能量需求对应等能线 $P\cdot t=E$，沿同一条线只能以功率换时间；整条线平移一格才叫效率改善，斜率不变。水平虚线是若干物理上限共同设定的功率天花板（其来源已在图内标出），它与等能线的交点给出最短时间 $t_{\min}$，这才是时间压不下去的原因。}
\label{fig:power-time}
\end{figure}


\subhead{四层阻抗}

第一层阻抗是功率。现实改变需要做功，搬动物体、驱动设备、压缩流程、提高产量、加快运输，本质上都要在单位时间内完成更多状态转移。能量除以时间就是功率：
\begin{equation}
P=\frac{\Delta W}{\Delta t},\qquad \Delta W\ \text{固定时}\ P\propto\frac{1}{t},
\label{eq:power}
\end{equation}
其中，$P$为功率，单位瓦特，即焦耳每秒；$\Delta W$为完成这次状态转移所需的功，单位焦耳；$\Delta t$为所用时间，单位秒。它成立的条件是功总量不因加快而减少，也就是路径、工艺和材料都没有变化，只是把同一件事做得更快；如果加快同时改变了工艺并降低了$\Delta W$，这个反比关系就要重新计算。由这个关系可以直接读出结论：时间减半意味着功率加倍，时间压缩一个数量级意味着功率上升一个数量级。相同能量消耗被压缩到更短时间，就要求更高功率输入，系统如果没有对应的供能、传输和控制能力，加速就会表现为过载。工业生产最容易看见这一点，产线提速不只是机器转得更快，还意味着供料、加工、检测、包装、物流和维护都要同步承受更高节奏，任何一个环节的功率截面不足，都会成为瓶颈，一个环节被强行加速，另一个环节就会堆积，速度不是在单点上发生的，而是在整条链上结算。

电力系统把这条反比关系写进了设备选型。要把功率的调整做得更快，靠的不是把同一台机组开得更猛，而是换成爬坡能力不同的设备：燃煤机组的出力调整通常在每分钟百分之几额定容量的量级，燃气机组明显更快，电池储能的响应落在毫秒到秒的量级，而它们的单位容量成本依次上升。一个电网能够承受的风光波动幅度，最终由这套设备的爬坡能力决定。同样的道理出现在制造与运输上，把交付周期从六周压到一周，通常不是让同一批设备连轴转，而是把一部分零件改成空运、把一部分工序改成并行、把一部分库存前置，省下的时间用更高的单位成本买回来。时间与功率之间的这条交换关系是连续的，压缩越到后面，需要的功率裕度越不成比例。

功率之外，第二层阻抗是热。任何真实过程都有损耗，计算有热，机械运动有摩擦，电力传输有损耗，化学反应有副产物，生产过程有废料，速度越高，单位时间内排出的废热越多。散热不是附属问题，而是加速的核心边界之一，很多系统不是不能更快计算、更快传输、更快制造，而是不能把热、损耗和副产物以同样速度排出去。数据中心、芯片、发动机、电池、工厂和城市交通都受这一点限制：能量可以更快注入，但废热也必须更快离开，输入通量提高，而排出通量不足，系统就会升温、降频、老化或失效。热告诉我们，加速不是只看能量能不能进来，还要看无用能量能不能出去。

功率和热之外，第三层阻抗是应力。物理对象不是理想符号，材料承受力，结构承受振动，设备承受磨损，连接件承受冲击，供应链承受波动，一个系统被更快推动时，内部薄弱点会更快暴露。机器加速会增加摩擦和疲劳，施工加速会增加质量控制压力，物流加速会减少缓冲空间，组织加速会减少沟通和修复时间。应力的危险在于，它往往不立即表现为失败，而是以微损伤形式累积，疲劳、磨损、松动、裂纹、质量波动，不一定来自一次大冲击，而来自高频负荷下的长期运行，加速让系统更快接近极限，也更快积累损伤。

还剩下最容易被忽略的一层，第四层阻抗是检测与修复。越快的系统，越需要更快发现错误：慢系统中，一个错误可能有时间暴露、定位、隔离和修复；快系统中，一个小错误可能在被发现前已经复制、扩散和放大。产线速度越快，检测窗口越短；金融交易越快，错误订单传播越快；内容生成越快，审核压力越集中；科研候选越多，验证瓶颈越重。检测和修复之所以是独立的截面，是因为它们的成本不随速度线性下降：把检测频率提高一倍，往往意味着设备、人员、标准和组织流程都要同步加密，而修复能力受制于现实中的定位时间和返工周期。

四层相加之后，得到一个与直觉相反的结论：快不一定是效率。效率是以较低总代价完成有效状态转移，速度只是单位时间内完成状态转移的程度，一个系统可以很快，但能耗高、热负荷大、损耗大、错误扩散快、维护困难；也可以较慢，但稳定、低耗、可修复。真正的问题不是能不能更快，而是更快之后，系统总账是否下降。

\subhead{速度与效率}

加速试验是工程上把时间换掉的正式做法，它的边界说明了这种交换能走多远。材料与器件的寿命评估普遍采用提高温度的办法加快反应速率，再按阿伦尼乌斯关系外推常温下的使用寿命，一项原本要跑几年的试验因此可以在几个月内完成。这类方法有一个硬条件：外推必须落在同一个失效机理之内，温度高到一定程度就会出现新的反应路径、新的应力状态或新的失效模式，此时外推出来的寿命不再是原来那个过程的寿命，而是另一个过程的寿命。试验的时间被压缩了，被试验的过程并没有被压缩，两者只在机理不变的区间里等价。这一类条件在加速的每一个应用场景里都会重演，凡是把时间参数直接调小的做法，都要先证明调小之后走的还是同一条路径。

混凝土的养护给出了同一条关系的正面用法。自然养护下强度按 28 天龄期发展，若把构件放进蒸压釜，在高压饱和蒸汽中养护十几个小时，就能达到接近自然养护 28 天的强度，预制管桩和高强构件走的正是这条路。省下的二十多天不是凭空消失的，它变成了蒸压釜、锅炉、耐压模具和持续供汽的能耗，并且只适用于能在密闭容器中处理的构件尺寸。同一份混凝土，用时间还是用能量把强度换出来，取决于现场更缺哪一个。

AI在这里尤其容易制造错觉。AI极大降低了搜索、表达、生成和初步筛选的时间，这种速度真实存在，也很有价值，它可以减少无效等待，减少盲目试错，缩小候选范围，提高早期判断效率。但当 AI 输出进入现实，它面对的不是同样可压缩的信息空间，而是材料、设备、身体、工艺、实验、组织和责任，前端生成越快，后端执行、验证和维护越可能成为通量瓶颈。把一万个候选筛到一百个可以很快，但把一百个候选变成可靠结论，需要实验、测试、重复、环境暴露和长期反馈，生成方案可以高频，现实验证不能无限高频，前端越快，后端越需要足够的验证截面，否则速度只是把压力从生成端转移到验证端。

后端自动化可以扩大并行度、减少等待，也可以通过改进路线减少步骤。剩下的限制要具体说明：哪些步骤确实必须依次完成，每步在当前条件下至少需要多久，哪些后果只有经过暴露才可观察。只有在这些条件固定时，链条才有相应的历时下限。

把加速的账本摊开看，收益与代价并不落在同一本账上。一个车间把节拍从 60 秒压到 45 秒，产出在纸面上提高了三分之一，但上料、换型、检测、包装和物流必须按新节拍重排，任何一环跟不上，多出来的半成品就会堆在工序之间，占用场地、资金和后续的返工时间。设备那一侧的账更直白：提高转速意味着更高的磨损速率和更短的检修间隔，装置利用率上升的部分，要从停机检修的时间里扣回来。速度在一个环节被兑现，在另一个环节被没收，加速不是把成本消灭，而是把成本从时间账挪到设备账和维护账上。

所以，加速不是一条单向通往效率的道路。真正有效的加速，是减少无效等待、减少返工、减少无意义摩擦，让系统以更低总代价完成同样状态转移；无效加速，则只是提高通量密度，把功率、热、应力、错误和维护压到后台。成熟系统不应盲目追求最短时间，而应寻找可承受速度区间，这个区间由供能、散热、承载、检测、修复和责任共同决定：低于这个区间，系统浪费时间；高于这个区间，系统消耗结构。真正的效率，不是把时间压到最短，而是在通量截面可以承受的范围内完成状态转移。加速的阻抗，就是时间压缩后的通量账本。时间不能被简单删除，它要么被浪费，要么被用于缓冲、散热、检测、修复和显影，AI可以帮助我们减少浪费的时间，但不能免费取消承载过程所需的时间，信息可以先到，现实还必须穿过通量截面。物理世界不是拒绝速度，它只是要求速度被系统承受。

以上讨论的都是同一种加速：把一个过程压进更短的时间窗口。还有另一种加速没有被触及，它不是压缩某个过程的持续时间，而是提高对同一个系统施加更新的频率。这两件事受不同的条件限制：前者取决于通量截面是否够大，后者取决于系统完成一次完整更新需要多久。如果一个输入的间隔短于系统吸收它的时间，系统就不是在加速，而是在被打断。

\subsection{计算不可约性：物理演化的步骤刚性}

时间刚性不只是过程不能任意压缩。它还有另一种更隐蔽的形式：现实系统不能被任意频率更新。AI很容易制造一种错觉，只要模型能够更快生成、预测和调度，现实系统也应当随之更快变化：模型可以每秒输出新判断，平台可以实时刷新数据，传感器可以连续采集信号，组织可以随时更新指标，算法可以不断重算方案，于是速度被误认为智能本身，实时被误认为效率本身。但现实系统不是屏幕，不是任何输入一到就完成状态改变，一个真实系统接受变化之后，还要经历传递、执行、反馈、修正和重新稳定，只有这些过程完成，系统才真正完成了一次更新。这就是时间的带宽。

\subhead{带宽与时间尺度}

对象的响应带宽说明输入变化有多少能够在输出中体现。材料、执行器、采购与组织各有动态过程，高频信号可能被衰减，也可能通过控制结构引发不同响应。新输入早于上一轮完成，并不自动构成失稳；需要检查输入、状态和闭环如何叠加。
\begin{equation}
\frac{\mathrm{d}x}{\mathrm{d}t}=\frac{u-x}{\tau},
\label{eq:relaxation}
\end{equation}
$x$和$u$采用相同状态单位，$\tau>0$为时间常数。输入固定时，偏差按指数衰减，一个时间常数后约剩37\%。输入持续变化时，系统表现为稳定的一阶低通，响应有幅值衰减与相位滞后。闭环振荡需要加入控制器、延迟等结构后另作分析。
\begin{equation}
f_c=\frac{1}{2\pi\tau},
\label{eq:cutoff}
\end{equation}
$f_c$为一阶系统的截止频率，单位为赫兹。正弦输入频率为$f$时，稳态幅值比为$[1+(2\pi f\tau)^2]^{-1/2}$；$f=f_c$时降到低频值的$1/\sqrt2$。高频响应因此受到衰减，但这个方程本身仍是稳定的一阶跟随系统。需要更多跟随带宽时，要改变对象或控制结构。

采样频率与对象跟随输入的能力是两个设计问题。较快采样可以提供更多状态信息，控制器还需结合延迟、噪声、增益和执行器限制来选择更新方式。一阶低通系统对高频输入的响应幅度减小；闭环是否振荡，则要检查控制器与对象共同形成的环路。采样快本身既不保证稳定，也不必然引起失稳。

部分过程可以通过更快的执行器、较短的协调链和更高的传感带宽缩短响应时间；另一些过程要改变材料或结构，才能改变其时间尺度。保护动作还要区分测量、判断与执行：快速测量配合滤波和阈值，可以同时保持响应与抗扰。免疫调节、断路保护和质量检查各有机制，所需节奏须按被识别的信号与错误后果设计。

因此，现实不是低速系统，而是多时间尺度系统。越接近接触、保护、平衡、故障隔离和瞬时安全，系统越需要快速反应；越接近目标、意义、制度、信任和长期结构，系统越需要保持稳定。中间还有计划、调度、库存、维护、协调和资源配置，它们既不能太慢，也不能被高频信号直接驱动。成熟系统不是让一个中心智能以同一种频率控制全部，而是在不同层级之间分配更新时间。

电力系统的分层是这一原则的成熟形态。故障发生后的最初几十毫秒由继电保护处理，它只做一件事：把故障设备从网络中切除，动作以毫秒计；随后的秒级到分钟级由安全稳定控制与调度处理，调整出力、切负荷、隔离区域，允许的时间尺度长得多；再往上是日前与中长期的市场、检修和规划，周期以天、月、年计。每一层的权限、动作速度和判断依据都不相同，越靠近物理极限的层越不允许人参与，越靠近目标的层越不允许自动接管。这套分层不是保守的产物，而是反复事故之后被写进导则的结构，其存在的理由正是任何单一时间尺度都无法同时应付故障和规划。

具身智能最能说明这一点。机器人不能把大模型直接接到每个电机上，身体不是屏幕，动作不是文字。语言理解可以较慢，视觉更新需要更快，运动规划需要把目标转化为轨迹，底层控制必须高速处理平衡、力、摩擦、滑移和误差。一个系统可以慢慢判断“这是杯子”，但手指接触杯子的瞬间，必须快速修正压力和姿态，语义正确不能替代力反馈。高层智能负责意义和目标，低层回路负责稳定和安全，中间必须有时间尺度转换，否则智能越强，身体越乱。电网也不是一个可以被任意实时重写的对象。故障保护必须极快，频率调节要快，负荷预测可以稍慢，市场调度更慢，电源结构和电网规划更慢。AI可以高频预测负荷，却不能高频重写整个电网，保护太慢会造成事故，规划太快会制造混乱，电网稳定依靠的不是无限实时，而是不同层级的更新时间各在其位。

\subhead{频率错配的形态}

供应链同样如此。需求信号可以实时跳动，产能、采购、库存、运输和质检却有周期，高频需求如果不加过滤地进入生产端，系统不会自动变得敏捷，反而会把短期波动转化为库存积压、缺货、返工和排程内耗。一个销售端的微小波动，在采购、生产和物流层层传递后，可能被放大成整条链的震荡，问题不在于需求信号太多，而在于慢系统被迫响应了不该响应的快波动。组织系统更典型：数据可以实时刷新，指标可以天天更新，AI可以不断提出建议，但人的信任、能力、协作习惯和责任关系不能每天重写。一个团队需要时间理解目标，需要时间形成节奏，需要时间把责任关系稳定下来，如果每天都有新重点，每周都有新目标，每次模型更新都有新建议，组织表面上更敏捷，实际却可能失去连续性，注意力被打碎，责任边界被冲淡，执行者不断切换方向，系统能量消耗在响应变化上，而不是完成任务上。

高频智能的真正危险，不是快，而是未完成即更新。上一轮变化还没有被系统吸收，下一轮变化已经到来；上一个指标还没有形成行为反馈，下一个指标已经替换；上一套方案还没有暴露后果，下一套方案已经下发。此时系统不再积累能力，而是在未结算的变化之间来回切换，它看起来一直在更新，实际上没有完成更新；看起来一直在响应，实际上没有形成稳定控制。频率错配通常有几种形态：第一是噪声化，输入太密，系统来不及区分真实变化和短期扰动，更多信息没有带来更多行动，只带来更多筛选、解释和协调；第二是打断化，高频建议不断切入低频流程，执行者不断重新理解、重新对齐、重新排序，系统的时间被切成无法完成工作的碎片；第三是振荡化，反馈存在延迟，调节却过于频繁，系统总是在追逐过去的误差；第四是责任漂移，自动建议越来越快，确认、追踪和承责却跟不上，失败后很难说清哪一次输出、哪一个接口、哪一个人或组织真正改变了系统状态。

把这四种形态放在一起，容易得出一个只往一个方向看的结论，似乎频率错配总是快扰动冲击慢系统。反方向的不匹配同样存在，而且后果更直接：需要快速响应的层级一旦慢下来，得到的不是低效，而是事故。故障保护迟了会烧毁设备，避障迟了会撞上行人，临床抢救迟了会失去窗口，反应时间在这里不是可以交换的成本，而是功能本身。所以时间尺度的问题从来不是快与慢哪一边更好，而是每一层的时间常数是否与它要处理的事情匹配，错配在两个方向上都会造成损失，只是损失的形式不同。

\subhead{更新时间接口}

因此，高频智能进入现实，不能只靠更快输出，而要靠更新时间接口。这个接口至少包括四类结构：限速、缓冲、保持和分层。限速不是压制智能，而是防止低频结构被高频输入反复重写；缓冲不是拖延，而是让短期波动不直接冲击慢系统；保持不是僵化，而是让一个决策有足够时间完成执行、反馈和修正；分层不是割裂，而是让快反馈留在局部，让慢判断留在结构层，让不同时间尺度不互相破坏。责任也是时间接口的一部分，现实系统不是纯粹信号系统，一个高频建议能否自动执行，不只取决于系统能否响应，还取决于失败后由谁承受。没有责任接口，高频智能越强，越可能把决策推向无人完整承接的自动过程，系统越实时，越需要清楚哪些变化可以自动执行，哪些必须等待确认，哪些只能作为提示，哪些不得进入行动。以此类推，“端到端”不能被简单理解为一个模型控制全部：端到端可以减少人工规则，也可以学习复杂映射，但现实执行仍然需要层级、限速和责任边界。真正的风险不只是黑箱，而是时间接口被黑箱化，一旦看不见哪些层在快，哪些层在慢，哪些反馈被过滤，哪些指令被保持，哪些失败由谁承受，系统就可能在表面一体化中失去稳定结构。

接口的建立还有一层容易被弄反的地方，它常被理解成给系统装一个减速器，仿佛把更新频率统一压低就安全了。实际要求不是慢，而是分层匹配：该快的层必须足够快，该慢的层不允许被催。AI的价值不在于让现实跟上模型速度。它可以高频观察，但不必高频干预；可以持续预测，但不必持续改变；可以快速生成方案，但必须稳定下发；可以增强决策，但不能替代现实系统自己的更新时间。成熟智能不是把世界推到模型频率，而是把模型能力嵌入世界的时间结构。实时也不是越快越好，而是在正确层级上及时：故障保护需要实时，战略规划不需要实时；接触反馈需要实时，组织文化不需要实时；交通避障需要实时，城市更新不需要实时。把不该实时的东西实时化，不是进步，而是把低频系统暴露在高频扰动中。

时间的带宽提醒我们：现实不是因为缺少智能才慢，而是任何持续系统都必须在变化和稳定之间分配时间。一个变化必须被传递、吸收、修正和稳定，才算真正完成，智能越快，越不能把输出频率误认为现实更新频率。高频智能不能直接等于高频现实，它必须知道哪些层可以快，哪些层必须慢，哪些变化必须被限速、缓冲、保持和分层。世界不是不能被智能驱动，而是不能被同一种时间尺度驱动。

接口能把更新频率安排到各层可以承受的区间，它改变不了另一件事：有些结果只有在过程走完之后才存在。对这类问题，提高刷新率不起作用，因为被等待的不是一个已经写好的答案，而是答案的形成。频率的边界之外，还立着另一层属于过程本身的边界。

\subsection{生命系统的时序依赖}

限速和缓冲处理的是更新频率，它们保证系统不被过快的变化打断，却没有解释另一种时间为什么会显得不可压缩。有些变化之所以不能被跳过，不是因为系统来不及吸收，而是因为它的结果在过程跑完之前并不存在。疲劳要在循环中累积，副作用要在代谢中显现，制度效果要在执行中暴露，生态后果要在网络中传播，这些事实不是预先存放在起点的答案，等待被更快地读出来，它们是过程本身的产物。

\subhead{材料、工程与生命}

材料世界最清楚地说明这一点。一个材料在实验室初测中表现良好，不等于它能在长期使用中保持稳定，疲劳需要循环，应力需要积累，腐蚀需要暴露，老化需要时间。很多失败不是在第一天出现，而是在长期微损伤中逐步形成，短期测试可以筛掉明显错误，却不能完全替代长期服役，时间不是等材料给出答案，而是材料在负荷、温度、水分、氧化、振动和维护条件中把隐藏缺陷展现出来的过程。石棉提供了一个极端的尺度：从吸入纤维到出现间皮瘤或肺纤维化，潜伏期通常在二十到四十年，任何把暴露与结果放在同一年里比较的评估方法都会给出错误的答案。

工程系统也一样。一个设计可以通过仿真，一个设备可以通过初测，一个系统可以在理想环境中运行，但可靠性要在噪声、故障、维护、规模化、边缘场景和用户行为中形成。很多问题不是设计者完全没有想到，而是只有当系统进入长期运行、被不同人使用、承受不同负载时才出现，可靠性不是图纸属性，而是过程属性。一个系统是否可靠，不只取决于它如何被设计，还取决于它如何老化、如何被维护、如何被误用、如何在故障中退化。检验环节安排了多长时间，本身就说明了这件事：半导体器件的寿命试验要在高温偏置下连续运行上千小时，全尺寸结构要做完设计寿命对应的循环次数，整车道路试验要跑完十几万公里，这些数字不是流程上的保守，而是让隐藏缺陷显形所必需的过程长度。

生命不是代码运行，而是物质过程按时序展开：基因表达有先后，蛋白合成、折叠、修饰、转运和降解都有速率，细胞分裂有周期，免疫反应有启动、放大、调节和消退，组织修复有阶段，发育和衰老本身就是长期过程。生命系统的真实不在序列本身，而在序列被细胞、组织、环境和时间逐步解释的过程中。这就是为什么药物研发不能因为AI加速分子发现而取消临床时间。AI可以更快找到候选分子，更快预测结构，更快筛选靶点，但候选物进入身体后，仍要面对吸收、分布、代谢、排泄、毒性、免疫反应、剂量、个体差异和长期反馈，短期有效不等于长期安全，局部靶点正确不等于全身网络可承受，动物模型有效不等于人群长期有效，生命系统不是静态容器，而是会响应、代偿、反弹和演化的复杂过程。这里的时间不是监管拖延，而是生物真实性的一部分：人体需要时间展现副作用，疾病需要时间展现进程，免疫系统需要时间展现反应，群体差异需要时间展现分布，一个药物真正进入现实，不是因为它在模型中看起来合理，而是因为它经受了生命过程的检验。

药物安全需要让作用在合适的人群、观察期与测量中显现。rentosertib在2026年9月到达三期首位患者给药阶段，是发现与开发链条上的新进展。\cite{insilico2026phase3} 它也显示，候选生成、进入临床与建立长期获益证据是不同里程碑。试验设计和监管能力可以改进，尚未发生的随访结果仍不能提前填入。

需要补一句的是，生命系统的时间并非全部不可压缩。筛查可以提前，生物标志物可以让损伤更早显形，适应性试验设计可以减少无效入组，剂量爬坡可以从更少的受试者中得到更可靠的起点，这些改进真实地缩短了个体差异被识别出来的时间。但它们压缩的是观测和筛选，不是暴露本身：药物在体内的代谢、组织在长期负荷下的累积、群体在真实使用中的分布，仍然要走完各自的步数。把可压缩的观测环节和不可压缩的暴露环节分开，才能判断一项加速究竟缩短了进程，还是只是把观测提前到了终点之前。验证不是附加环节，而是生命系统把候选转化为事实的方式。

\subhead{组织与社会}

组织和社会系统也有时间证明。一个制度设计可以在纸面上合理，一个政策模型可以在模拟中有效，一个管理方案可以在逻辑上清楚，但它进入组织后，会遇到激励、信任、习惯、解释、抵抗和套利，人不会像变量一样被直接赋值。制度效果要在执行中展现，信任要在重复互动中形成，责任要在事件中被确认，组织能力要在长期协作中积累，社会事实不是方案的静态推论，而是人群反应在时间中形成的结果。很多改革初期看起来成功，后来才暴露代价；很多平台机制一开始提高效率，后来才出现滥用、内卷和系统脆弱；很多教育、医疗、城市治理方案短期指标改善，长期后果却复杂得多。一个制度不是因为逻辑完整就已经成立，而是要看它如何被人理解、利用、绕行和反制，社会系统会学习，也会变形；会适应，也会反作用，时间不是外在等待，而是社会后果形成的空间。

生态系统则把这种结构推到更长尺度。一个化学品、一个外来物种、一种土地利用方式、一项水利工程，短期可以有明确收益，长期却可能通过食物链、水循环、土壤、物种关系和气候反馈展现不同后果。生态不是线性机器，而是多重反馈网络，局部改变可以沿关系链传播，短期收益可以在长期中转化为系统代价，生态后果之所以难以完全提前证明，不是因为人不够聪明，而是因为后果要在开放网络中运行后才形成。这类过程的时间尺度往往以十年计：表层土壤的形成以每厘米数百年计，被采伐的森林恢复到原有结构与物种组成通常需要几十年以上，湖泊在营养盐停止输入之后仍会在沉积物中缓慢释放磷，水质改善滞后于减排若干年。想在这些系统里提前拿到结论，需要的是观测序列，不是更快的推理。

这些例子背后有几种更基本的过程。第一是累积，微损伤、毒性、信任损耗、生态压力，不一定一次出现，而是在重复循环中增加。第二是阈值，系统可以长期看似稳定，直到某个边界被越过，失效突然发生。第三是适应，生命、组织、市场和制度会对干预作出代偿、绕行和反制。第四是扩散，局部后果会沿材料裂纹、代谢网络、社会关系和生态链条传播。时间证明不是单一等待，而是这些过程把不可见状态变成可观察事实。

不过，社会过程和材料过程之间有一处关键差别，它使前者比后者更难收敛。材料在负荷下如何失效，不会因为被观测而改变；人和组织会。政策一旦公布，预期就重新形成；指标一旦进入考核，行为就围绕指标重组；监督一旦加强，被监督的部分就向未被监督的部分转移。于是同一个干预在不同时间点测得的效应并不稳定，过程不是一条固定的轨迹，而是随观测和规则一起变化的轨迹。正因如此，社会领域里“再跑一遍就知道结果”的做法常常失效：第一次运行的结果已经改变了第二次运行的初始条件。社会事实因此既需要时间才能形成，又不能在形成之后被当作一次性的定论。

\subsection{时间的证明：过程作为信息的来源}

智能时代最容易被低估的，不是生成能力，而是某些事实形成所需要的时间。模型可以快速生成答案、设计、分子、策略和预测，它可以在表示空间中压缩搜索，把许多过去需要长期尝试的可能性提前呈现出来。但现实问题不只是在可能性空间中找到一个看似合理的解，很多真实结果不是预先完整藏在起点里的答案，而是在材料、生命、组织和生态的运行过程中逐步形成、暴露并固定下来，它们不能被一次推理完全替代，也不能被一个静态指标直接证明。这就是时间的证明。

\subhead{过程作为信息的来源}

时间证明不是简单等待，也不是对智能不足的补偿。它是让候选进入过程，让系统在循环、暴露、代谢、使用、扩散和反作用中把隐藏后果变成事实。一个候选方案在表示空间中可以已经完整，但它在现实中是否成立，要看它进入材料后如何承受循环，进入身体后如何被代谢，进入组织后如何被理解和反制，进入生态后如何沿网络扩散，时间不是验证之外的空耗，而是某些事实形成和展现的条件。这可以借用计算不可约作为启发，但不能把它粗暴理解为“复杂系统都不可预测”。很多系统可以被模拟，可以被粗粒度预测，可以用代理指标提前判断方向，真正关键的是区分两类东西：一类可以通过计算、推理和仿真提前压缩，另一类必须由过程运行来给出。

所谓后者，指的是那些结果由过程本身携带的东西：材料的疲劳史、人体的代谢反应、组织的适应路径、生态的长期反馈，都不是静态候选的简单标签，而是系统在时间中与环境相互作用后形成的事实。这类过程的时长有一个下界：
\begin{equation}
T\;\ge\;N\,\delta t,
\label{eq:steps}
\end{equation}
其中$T$为过程完成所需的总时间，单位秒；$N$为该过程必须依次完成的演化步数，无量纲；$\delta t$为单步所需的最短时间，单位秒，由该步对应的物理过程决定，例如一次疲劳循环、一轮代谢、一个繁殖世代或一段暴露期。两侧量纲一致，都是时间，这个关系成立的条件是步与步之间存在依赖：后一步的初始状态由前一步的结果决定，因此不能并行，也不能跳过。如果过程可以被分解为互不依赖的子过程，$N$可以被并行度分摊，这个下界就不再收紧。计算不可约的含义正在这里：对这类过程，不存在一个更聪明的算法把$N$压缩成远小于原值的数，你只能把同样的步数走完，能改变的只是每一步的$\delta t$，而$\delta t$由物理过程的速率常数决定，不由推理速度决定。

下界的形式简单，但它的两个量都不容易确定，判断一项加速是否真的缩短了过程，难处正在这里。$N$不是过程本身的固有属性，它取决于把什么算作一步：材料疲劳可以按一次应力循环计，也可以按一段裂纹扩展的阶段计，切分尺度不同，$N$就差出几个数量级，而粗粒化不能改变$T$本身，只能改变我们数出来的步数。$\delta t$也不总是可以由参数直接读出，同一轮代谢在不同个体身上快慢不同，同一段暴露期在不同气候和负荷下长短不同。把这两个量分开看，可以得到一条实用的判据：压缩$\delta t$的技术进步是真实的，它改的是过程速率；把$N$改小的说法需要额外证明，因为只有在依赖结构真的被拆开时，$N$才会下降，否则减少的只是被记录下来的步数，过程本身仍然要跑完。

不过，这条区分有一种对称的误用。既然有些事实只能由过程给出，就容易被推广成一切长期判断都不必提前做，一切预测都不可信，只能等结果自己出现。这同样不成立，因为过程不可跳过，不等于过程不能被打断、不能被提前识别、不能在中途被重新定向：失效的早期信号、代谢的中间产物、组织的态度变化、生态的指示物种，都是在过程尚未跑完时就出现的信息，它们不给出终局结论，却能把验证从盲目等待变成有方向的观测。真正需要区分的是终局事实和过程信号：前者必须走完全程，后者可以在中途采集，把两者混在一起，就会在应该继续观测的时候提前下结论，或者在应该等完的时候反复要求结论。

过程信号的价值可以从几个具体指标上看到。糖化血红蛋白反映的是红细胞寿命期内、也就是约两到三个月的平均血糖水平，单次血糖测量只给出一个时点，它给出的是这段时间里累积暴露的结果，因此它比瞬时指标更适合判断长期控制状况，但也因此无法回答今天早上发生了什么。金属结构内部的裂纹在扩展到临界尺寸之前会发出声发射信号，监测这些信号可以在结构失效之前给出方向，却不能在第一步循环里就给出寿命。这类信号都遵守同一条分工：它们把不可见过程的一部分提前暴露出来，让观测有方向，但不替代过程本身走完。

\subhead{可压缩与不可压缩}

AI的作用仍然非常重要。它可以减少候选空间，发现更可能的路径，提前识别风险，设计更好的实验，模拟更多场景，寻找早期信号，压缩无效试错，它能让验证更聪明、更有方向、更少浪费。但AI不能取消现实验证本身，它能减少要等待的对象，不能取消等待的必要性；能提高预测能力，不能让过程完全跳步；能提前发现许多错误，不能让世界无需亲自运行。因此，成熟判断不在于否认AI的速度，而在于区分两种时间：搜索时间可以被压缩，候选生成时间可以被压缩，形式推演时间可以被压缩，但事实形成的时间不能被随意取消。把可以计算的部分拖慢，是低效；把必须由过程给出的部分跳过，是危险。真正的技术进步，不是把所有时间都视为浪费，而是识别哪些时间是无效等待，哪些时间是现实给出事实的必要过程。

这也是为什么生成越便宜，时间证明反而越重要。一个模型可以生成一千个药物候选、工程方案、商业策略或政策文本，但现实不能同等速度扩大临床、长期试验、工程服役、组织反馈和生态观察，信息空间膨胀之后，稀缺不再只是想法，而是能够让少数关键候选穿过真实过程的验证能力。AI降低的是提出可能性的成本，时间证明承担的是把可能性转化为可信事实的成本。有些过程不能被一个简单捷径完整压缩，因为过程本身会产生新的约束、新的状态和新的反馈：你可以预测轨迹，但轨迹仍要被走过；你可以模拟风险，但风险仍要被现实检验；你可以压缩搜索，但不能取消过程提供的信息。现实不是一张可以提前填完的表，而是一条必须展开的过程。

结果是否能被快速求出，取决于问题结构、算法、所需精度与信息来源。初始条件完整，不保证存在快捷算法；过程逐步展开，也不保证任何近似或重设计都无效。只有在必须经历的顺序步骤、单步条件与观察要求得到说明后，才能把$T\geq N\delta t$用于给定过程。计算不可约性不替每一项工程等待提供证明。

这不是技术悲观主义。真正成熟的技术应当尊重时间证明，并让它更有效率：更好的传感器可以更早发现疲劳，更好的生物标志物可以缩短临床判断，更好的数字孪生可以减少实体试错，更好的统计设计可以提高实验效率，更好的AI可以帮助选择最值得验证的路径。技术的任务不是否认时间，而是减少无效等待，放大有效反馈。关键在于边界：哪些结果可以由计算提前给出，哪些结果只能由过程给出；哪些风险可以通过模型排除，哪些风险必须由暴露、使用和长期反馈展现；哪些候选已经足以进入下一步，哪些候选必须停在表示空间。智能时代最重要的判断力之一，就是区分“可以提前压缩的时间”和“必须保留的时间”。时间的证明提醒我们：真实不是所有时候都以答案形式预先存在，有些真实要经过疲劳、代谢、免疫反应、组织适应、生态反馈和长期使用，才会形成可依赖事实，它们不是没有被聪明人想到，而是必须由世界运行出来。AI可以压缩寻找答案的时间，但不能取消答案成为事实的时间，有些事实，只能由过程给出。

过程给出事实之后，还有一道关口。一个候选通过了验证，只说明它在自己的尺度上成立，不说明它能嵌入既有的系统。现实系统由相互依赖的设备、人员、制度、习惯和利益构成，改动一个节点的实际代价，取决于这个节点在结构中的位置和它与多少其他部分相连。同样一项通过验证的技术，装在不同位置，遇到的阻力可以相差一个数量级。
